
\begin{remark}[Steered block-refined triangular solver; conditional; AI-checked, pending human review]
\label{rem:sde-block-steered}
This remark comes from an AI-only pass (\file{research/audits/reviews/2026-11-05-rees-sde-two-relation.md}), checked and corrected by a second, adversarial AI pass (\file{research/audits/reviews/2026-11-05-rees-sde-block-check.md}); no human has checked it, nothing else in the paper uses it, and it gives no worst-case improvement. Keep the notation of Proposition~\ref{prop:triangular-sde}, with $Q_\nu=g^{\nu}P$ and $Q_\nu^\circ=Q_\nu\cup\{0\}$. Let $S$ be a circuit modulo $U$ whose dependency $\delta$ has entries in at least two cosets, and split $S$ into blocks $G_1,\dots,G_E$, each contained in one set $\{i:\mathrm{cos}(\delta_i)=c_b\}$; put $W_b=\sum_{i\in G_b}\delta_i\lambda_i$. Then the classes of the $W_b$ modulo $U$ span a space $\Omega$ of dimension $E-1$, and their only relation is $\sum_bW_b=e:=\sum_{i\in S}\delta_i\lambda_i\in U$. For $\beta\in\Fp^E$ the weights $x_i=\delta_i(\beta_b+s)$, $i\in G_b$, lie in $P\cup\{0\}$ exactly when $s\in\mathrm{Good}(\beta)=\{s:\beta_b+s\in Q_{-c_b}^\circ\text{ for all }b\}$, and their value is $\sum_b\beta_bW_b+s\,e$; every $x\in(P\cup\{0\})^S$ whose value lies in $\Omega$ modulo $U$ has this form. By the Weil bound~\cite[Theorem~5.41]{LidlNiederreiter1997}, $|\mathrm{Good}(\beta)|\ge\kappa^{-E}\bigl(p-E-(\kappa^E-1)(E-1)\sqrt p\bigr)$ whenever the $\beta_b$ are pairwise distinct. Let $E_{\mathrm{cg}}(E,\kappa)$ be the set of primes $p\equiv1\pmod\kappa$ at which $\mathrm{Good}(\beta)=\emptyset$ for some $E'\le E$, some classes $c_1,\dots,c_{E'}$ not all equal and some pairwise distinct $\beta_1,\dots,\beta_{E'}$, and let $G_E(p)$ be the least $|\mathrm{Good}(\beta)|$ over such data. Exhaustive computations up to the Weil bound, repeated by the second pass with independent code (finite computations, AI-checked, pending human review), give $E_{\mathrm{cg}}(3,2)=\emptyset$ and $E_{\mathrm{cg}}(3,3)=\{7,13,19,31,37,43,61,67,73,79,97,103,109,127,139,151,157,163,181,193,211,241\}$. So for $p\notin E_{\mathrm{cg}}(E,\kappa)$ a circuit realises every point of $\Omega$ off the at most $\binom{E}{2}$ hyperplanes $\{\beta_b=\beta_{b'}\}$ with $c_b\ne c_{b'}$, by drawing $s$ uniformly until $E$ coset tests pass, about $p/G_E(p)$ draws; by part~(a) of Remark~\ref{rem:sde-subquadratic} it never realises a whole plane. Running the triangular recursion with these blocks, and drawing the free scalar $s$ of each round uniformly from $\mathrm{Good}(\beta)$ in the back substitution, moves the next residual along $e$ and steers it off the bad hyperplanes. Suppose that $E\ge\kappa$, that $p\notin E_{\mathrm{cg}}(E,\kappa)$, that the columns of every round, except a last round whose dependency lies in one coset, form a circuit modulo the current space, and that no residual offset difference $\beta_b-\beta_{b'}$ with $c_b\ne c_{b'}$ is identically zero as a linear function of the drawn scalars. Both conditions concern the forward pass, including its choice of blocks, and can be checked in polynomial time once that pass is fixed; when one fails the procedure may abort, and it never returns a wrong solution, since its output is checked. Under these conditions each back substitution returns a nonzero solution with probability at least $1-L\binom{E}{2}/G_E(p)$, where $L\le\lceil a/(E-1)\rceil$ is the number of rounds, and the recursion uses at most
\[
  M_E(a)=1+m(a+1)-\tfrac{(E-1)m(m-1)}{2}+[r>0](r+1),\qquad a=m(E-1)+r,\ 0\le r<E-1,
\]
columns, which is $(a+1)(a+2)/2$ for $E=2$, $\lfloor(a+2)^2/4\rfloor+[a\text{ odd}]$ for $E=3$, and $a^2/(2(E-1))+O(a)$ in general; the expected time is polynomial in $a$ and $\log p$ times $\kappa^E$. For i.i.d.\ uniform columns both conditions hold with probability at least $1-\bigl(M_E(a)+L\binom{E}{2}\bigr)/(p-1)$; this is a property of the columns, and repeating the back substitution does not reduce it. Hence, for i.i.d.\ uniform columns in $\Fp^a$ with $a\ge2$ and $E=3$, a randomised polynomial-time procedure with $\lfloor(a+2)^2/4\rfloor+[a\text{ odd}]$ columns fails with probability at most $1/16$ over the columns, and otherwise with probability at most $8^{-k}$ after $k$ back substitutions, for $\kappa=2$ at every prime $p\ge64a^2$ and for $\kappa=3$ at every prime $p\ge\max(2750,1000a^2)$ (the second range uses an exhaustive computation of $G_3(p)$ for $p\le11000$); with $E=\lceil\log_\kappa a\rceil$ and $p\ge\poly(a)$ it uses $O(a^2/\log a)$ columns. The steering condition is genuinely needed: there are column sets with $M_E(a)\ge\kappa a+1$ columns, hence solvable by the Chevalley--Warning theorem, on which it fails and every back substitution fails (planted instances, exact computations, AI-checked, pending human review: for $\kappa=3$, $E=3$ and $a\in\{8,10\}$, where the blocks are forced to be the coset classes, and for $\kappa=2$, $E=3$ and $a=6$ with a fixed block rule). So for arbitrary columns the best arity proved in this paper remains $(a+1)(a+2)/2$. The labels of Lemma~\ref{lem:s3-law} are i.i.d.\ uniform on the unknown subspace $K_d^\perp$, and the procedure commutes with invertible linear maps of the columns, so run on coordinates in the span of its bank, which is $K_d^\perp$ except with probability at most $(p^{a_d}-1)/((p-1)p^{M})$, it behaves as on i.i.d.\ uniform columns in $\Fp^{a_d}$; with the public arity $M_E(g_d)$ it would therefore give a polynomial-time option that uses the label law, but only at primes large compared with $g_d^2$ and with the number of nodes of the recursion. The query counts of Theorem~\ref{thm:main}, Remark~\ref{rem:combined-schedule}, Theorem~\ref{thm:free-coordinate} and the tables do not use it, and it would not change $E_W(c)$, because its prime condition grows with $g_d$.
\end{remark}
